\documentclass[10pt,a4paper]{beamer}

\usepackage[ngerman]{babel}
\usepackage[T1]{fontenc}
\usepackage[utf8]{inputenc}
\usepackage{varioref}
\usepackage{hyperref}
\usepackage{cleveref}
\usepackage{amsmath}
\usepackage{amsfonts}
\usepackage{amssymb}
\usepackage{makeidx}
\usepackage{graphicx}
\usepackage{csquotes}
\usepackage{listings}
\usepackage{color}
\usepackage{xcolor}
\usepackage[most]{tcolorbox}
\usepackage{amssymb}
\usepackage{lmodern}
\usepackage{verbatim}

% Umlaute und ß in Listings
\lstset{basicstyle=\ttfamily}
\lstset{literate=%
  {Ö}{{\"O}}1
  {Ä}{{\"A}}1
  {Ü}{{\"U}}1
  {ß}{{\ss}}1
  {ü}{{\"u}}1
  {ä}{{\"a}}1
  {ö}{{\"o}}1
}

% Farben für Listings
\definecolor{codegreen}{rgb}{0,0.6,0}
\definecolor{codegray}{rgb}{0.5,0.5,0.5}
\definecolor{codepurple}{rgb}{0.58,0,0.82}
\definecolor{backcolour}{rgb}{0.95,0.95,0.92}
 
\lstdefinestyle{mystyle}{
    backgroundcolor=\color{backcolour},   
    commentstyle=\color{codegreen},
    keywordstyle=\color{magenta},
    numberstyle=\tiny\color{codegray},
    stringstyle=\color{codepurple},
    basicstyle=\footnotesize,
    breakatwhitespace=false,         
    breaklines=true,                 
    captionpos=t,                    
    keepspaces=true,                 
    numbers=left,                    
    numbersep=5pt,                  
    showspaces=false,                
    showstringspaces=false,
    showtabs=false,                  
    tabsize=2,
	language=[LaTeX]{TeX}
}
 
\lstset{style=mystyle}

\definecolor{hbbkBlueish}{HTML}{5e7b87}
\definecolor{hbbkReal}{HTML}{1A4354}
\setbeamercolor{logo}{bg=white}  %controls the color of the logo area
\setbeamercolor{author in sidebar}{fg=white} 
\setbeamercolor{palette primary}{bg=hbbkBlueish,fg=white}
\setbeamercolor{palette secondary}{bg=hbbkBlueish,fg=white}
\setbeamercolor{palette tertiary}{bg=hbbkBlueish,fg=white}
\setbeamercolor{palette quaternary}{bg=hbbkBlueish,fg=white}
\setbeamercolor{structure}{fg=hbbkBlueish} % itemize, enumerate, etc
\setbeamercolor{section in toc}{fg=hbbkReal} % TOC sections


\usetheme{PaloAlto}
\renewcommand{\footnotesize}{\small}
\newcommand{\pftn}[1]{\let\thefootnote\relax\footnotetext{\tiny #1}}
\newcommand{\ftn}[2]{\footnote[#1]{\tiny #2}}
\newenvironment{hlbox}{\begin{tcolorbox}[enhanced,colback=white,colframe=white,sharpish corners,fuzzy halo=0.5mm with lightgray]}{\end{tcolorbox}}

\logo{\includegraphics[width=0.155\linewidth]{hbbk-logo}}
\author{Luca Kiebel}

\author{Luca Hülsmann \and Luca Kiebel \and Luca Hartmann}
\title{Big Data}
\date{\today}
\setlength{\itemsep}{10pt}
\begin{document}
\begin{frame}
\titlepage
\end{frame}

\begin{frame}
\frametitle{Inhaltsverzeichnis}\tableofcontents
\end{frame}

\section{Was ist Big Data?}

\subsection{Definition}
\begin{frame}
\frametitle{Definition}
\framesubtitle{Was ist Big Data?}
Ich bin cool
\end{frame}

\subsection{Wofür wird Big Data genutzt?}
\begin{frame}
\frametitle{Wofür wird Big Data genutzt?}
\framesubtitle{Was ist Big Data?}
\begin{itemize}
\item Verhaltensanalyse
	\begin{itemize}
	\item Erkennen was wann gekauft wird
	\item Erkennen wieso etwas gekauft wird
	\end{itemize} \pause
\item Targeted Advertising
	\begin{itemize}
	\item Personalisierte Werbung
	\item Coupons anbieten 
	\end{itemize} \pause
\item Predictive Support
	\begin{itemize}
	\item Erkennen von Potenziellen Schwachstellen
	\item Austauschen von Teilen bevor Abnutzung bemerkbar wird
	\end{itemize}
\pftn{https://www.slideshare.net/Dell/big-data-use-cases-36019892}
\end{itemize}
\end{frame}

\subsection{Wer nutzt Big Data?}
\begin{frame}
\frametitle{Wer nutzt Big Data?}
\framesubtitle{Was ist Big Data?}
\begin{itemize}
\item \( 1/3\) deutscher Unternehmen
\pause
\item 2014: 23\%
\pause
\item Hilfe zu relevaten Entscheidungen bei 80\%
\pause
\item 2020: 1,7 MB pro Person in einer Sekunde \ftn{1}{https://blog.capterra.com/10-surprising-big-data-statistics/}
\pftn{https://www.bitkom.org/Presse/Presseinformation/Jedes-dritte-Unternehmen-nutzt-Big-Data.html}
\end{itemize}
\end{frame}

\section{Welche Vorteile haben Unternehmen durch Big Data?}
\begin{frame}
\frametitle{Welche Vorteile haben Unternehmen durch Big Data?}
Inhalt...
\end{frame}

\section{Welche ethischen Probleme kommen durch Big Data auf?}
\begin{frame}
\begin{itemize}
\frametitle{Welche ethischen Probleme kommen durch Big Data auf?}
\item Eigentum von Daten ist in Deutschland nicht möglich (somit eine Rechtslücke)
\pause
\item Daten als solche unterliegen keinem Schutzrecht wie etwa dem Urheberrecht oder dem Patentschutz \ftn{https://www.computerwoche.de/a/wem-gehoeren-die-daten-im-internet-of-things,3328337}
\item nur ethisch fragwürdig, wenn der Kunde nichts davon weiss und keine Möglichkeit habe, allfällige Fehler zu korrigieren
\end{itemize}
\end{frame}

\section{Bewertung}
\begin{frame}{Bewertung}
Inhalt...
\end{frame}

\end{document}
